\documentclass[10pt,a4paper]{amsart}
\usepackage[margin=19mm]{geometry}
\usepackage{amsmath,amssymb,mathtools}
\usepackage[scr=rsfs,cal=euler]{mathalfa}
\usepackage{xfrac}
\usepackage[unicode,hidelinks,bookmarks=false]{hyperref}
\newcommand{\ie}{i.e.\ }
\newcommand{\cf}{cf.\ }
\newcommand{\resp}{respectively\ }
\newcommand{\Subsection}{\subsection}
\providecommand{\nobreakdash}{\nobreak\hbox{-}}
\setlength{\emergencystretch}{2em}
\multlinegap0pt
\usepackage{amssymb}
\usepackage{mathtools}
\usepackage{booktabs}
\usepackage{tikz}
\usepackage[shortlabels]{enumitem}
\newtheorem{lemma}{Lemma}
\newtheorem{theorem}{Theorem}
\newcommand{\ZZ}{\mathbb{Z}}
\newcommand{\colA}{\mathrm{A}}
\newcommand{\colB}{\mathrm{B}}
\newcommand{\colC}{\mathrm{C}}
\title{\textbf{Exact Dominion of the Prism Graph:\\ Enumeration by Congruence Class via Cyclic Words}}
\author{Julian Allagan}
\date{Source: arXiv:2601.03488v1 (CC BY 4.0)}
\begin{document}
\maketitle

\noindent\emph{Source and licence.} \textbf{Exact Dominion of the Prism Graph:\\ Enumeration by Congruence Class via Cyclic Words}. Version arXiv:2601.03488v1 (CC BY 4.0), distributed by arXiv under the Creative Commons Attribution 4.0 International licence, http://creativecommons.org/licenses/by/4.0/, as recorded in the arXiv licence metadata for this exact version and shown on the abstract page https://arxiv.org/abs/2601.03488v1. Full licence notice: \texttt{licenses/arxiv-2601.03488-CC-BY-4.0.txt}.\par\smallskip

\noindent\textit{Editorial scope.} Counted unit: the article's theorem thm:mod2-quadratic (source0.tex lines 591--593). The article's complete original proof of it is delivered with it (source0.tex lines 595--709, one \texttt{proof} environment of 115 lines). No mathematics is written, repaired or re-derived: the statement and the proof are the article's own lines, in the article's own notation, and no retained line is substituted or re-flowed.  Retained uncounted as the source-local context the proof needs: lines 512--515; lines 531--538; lines 576--579.  Nothing was omitted from any retained span, so the package carries no omission record.  No retained line contains a citation, so the wrapper carries no bibliography and no bibliography entry.  No retained line contains a figure environment, an image-inclusion command or drawing code, so no image is delivered and the package declares no asset dependency.  Editorial formatting only, in addition: an AMS \texttt{amsart} preamble that restates the article's own declarations and macros used by the retained lines, namely the environments \texttt{lemma}, \texttt{theorem} on separate counters, together with the standard packages those lines need.  The retained environments print in this document as Lemma 1, Theorem 1 in order of appearance, whereas the article prints the same items with the numbering of its own full document.  The complete native source of the article as distributed in the arXiv e-print for this exact version is bundled unchanged as \texttt{sources/192229/source0.tex}.
\begin{lemma}[Two-gap structure]\label{lem:mod2-two-gap}
Every minimum dominating word for $G_{4t+2}$ has exactly two size-$2$ gaps between
consecutive $\colC$'s, and all remaining gaps have size $1$.
\end{lemma}

\begin{lemma}[Monochromatic size-$2$ gaps]\label{lem:mod2-gap-mono}
Let $n=4t+2\ge10$, and let $w\in\{\colC,\colA,\colB\}^n$ be a minimum dominating word.
Then every size-$2$ gap has the form
\[
\colC\,\varepsilon\,\varepsilon\,\colC
\qquad\text{with }\varepsilon\in\{\colA,\colB\}.
\]
\end{lemma}

\begin{lemma}[Forced backbone segment]\label{lem:mod2-backbone}
Let $w$ be a minimum dominating word for $G_{4t+2}$.  Cutting the cycle immediately after
a size-$2$ gap, the next $4t$ letters form a backbone interval.
\end{lemma}

\begin{theorem}[Quadratic dominion for $n\equiv 2\pmod4$]\label{thm:mod2-quadratic}
If $n=4t+2\ge10$, then $\zeta(G_n)=n(n+2)$.
\end{theorem}

\begin{proof}
Let $\mathcal{L}_n$ denote the set of minimum dominating words in
$\{\colC,\colA,\colB\}^n$ for $G_n$, so $|\mathcal{L}_n|=\zeta(G_n)$.

For $w\in\mathcal{L}_n$, let $\mathrm{end}(w)\subseteq\ZZ_n$ be the set of indices $j$
for which
\[
(w_{j-4},w_{j-3},w_{j-2},w_{j-1})
=(\colC,\varepsilon,\varepsilon,\colC)
\quad\text{for some }\varepsilon\in\{\colA,\colB\}.
\]
By Lemmas~\ref{lem:mod2-two-gap} and~\ref{lem:mod2-gap-mono},
$|\mathrm{end}(w)|=2$ for every $w$.

Let $\rho$ be the cyclic shift $(\rho^k w)_i:=w_{i+k\ (\mathrm{mod}\ n)}$ and define the
anchored set
\[
\mathcal{A}_n:=\{w\in\mathcal{L}_n:\ 0\in\mathrm{end}(w)\}.
\]

\begin{lemma}[Trivial stabilizer]\label{lem:trivial-stab}
If $w\in\mathcal{L}_n$ and $\rho^k(w)=w$ for some $k\in\ZZ_n$, then $k\equiv 0\pmod n$.
\end{lemma}

\begin{proof}
Assume $\rho^k(w)=w$. Then $\mathrm{end}(\rho^k w)=\mathrm{end}(w)$, but by definition
$\mathrm{end}(\rho^k w)=\mathrm{end}(w)-k$ in $\ZZ_n$. Hence
\begin{equation}\label{eq:end-invariant}
\mathrm{end}(w)=\mathrm{end}(w)-k.
\end{equation}
Write $\mathrm{end}(w)=\{e_1,e_2\}$ with $e_1\neq e_2$ (recall $|\mathrm{end}(w)|=2$).
From \eqref{eq:end-invariant}, translation by $-k$ permutes $\{e_1,e_2\}$.
Thus either $k\equiv 0$ (the identity translation), or $-k$ swaps the two endpoints, i.e.,
\[
e_1-k\equiv e_2\quad\text{and}\quad e_2-k\equiv e_1 \pmod n,
\]
which implies $2k\equiv 0\pmod n$. Since $n=4t+2$ is even, this yields $k\equiv n/2\pmod n$
as the only remaining possibility.

Suppose $k\equiv n/2$. Then $\mathrm{end}(w)=\{a,a+n/2\}$ for some $a$.
Cut the cycle immediately after the size-$2$ gap ending at $a$.
By Lemma~\ref{lem:mod2-backbone}, the next $4t$ letters form a backbone interval, so the
\emph{other} size-$2$ gap must end exactly $4t$ steps later, i.e., at $a+4t$ modulo $n$.
Therefore
\[
a+4t \equiv a+\frac{n}{2}\pmod n.
\]
With $n=4t+2$, this becomes $4t\equiv 2t+1\pmod{4t+2}$, equivalently $2t\equiv 1\pmod{4t+2}$,
which is impossible for $t\ge 2$. Hence $k\not\equiv n/2$, and thus $k\equiv 0$.
\end{proof}

\begin{lemma}[Orbit--anchor relation]\label{lem:orbit-anchor}
Every rotation orbit in $\mathcal{L}_n$ has size $n$ and contains exactly two anchored
words.  Consequently,
\begin{equation}\label{eq:orbit-formula}
|\mathcal{L}_n|=\frac{n}{2}\,|\mathcal{A}_n|.
\end{equation}
\end{lemma}

\begin{proof}
Triviality of stabilizers follows from Lemma~\ref{lem:trivial-stab}, so each orbit has
size $n$.  If $\mathrm{end}(w)=\{e_1,e_2\}$, then precisely the shifts
$\rho^{e_1}(w)$ and $\rho^{e_2}(w)$ place an endpoint at $0$, and no other shift does.
\end{proof}

The second size-$2$ gap is determined by the location of the length-$4$ window
\[
(\colC,\varepsilon,\varepsilon,\colC)
\]
that ends at its rightmost $\colC$ (equivalently, by the index $j$ with
$(w_{j-3},w_{j-2},w_{j-1},w_j)=(\colC,\varepsilon,\varepsilon,\colC)$).
Besides the $n-1$ interior choices $j\in\{1,2,\dots,n-1\}$, there are exactly two
wrap-around choices, corresponding to windows that cross the cut between indices
$n-1$ and $0$:
\[
\text{(wrap 1)}\quad (w_{n-2},w_{n-1},w_0,w_1)=(\colC,\varepsilon,\varepsilon,\colC),
\]
\[
\text{(wrap 2)}\quad (w_{n-1},w_0,w_1,w_2)=(\colC,\varepsilon,\varepsilon,\colC).
\]
Thus the second gap has precisely $(n-1)+2=n+1$ admissible window-positions relative to
the anchor at $0$.

The second gap may be placed at any of the $n-1$ interior positions $\{1,2,\dots,n-1\}$
or may wrap around the cyclic boundary in two distinct ways, requiring positions beyond
the standard index range.
To parameterize all possible placements of this second gap, including the two ways it
may straddle the linear boundary, introduce the extended index set
\[
\widetilde{\ZZ}_n:=\{1,2,\dots,n-1\}\cup\{n,n+1\}.
\]
Each element of $\widetilde{\ZZ}_n$ encodes a unique length-$4$ window where a second
monochromatic gap of type $\varepsilon$ may be implanted.  Conversely, every anchored
minimum word arises uniquely from such a choice, by Lemmas~\ref{lem:mod2-gap-mono}
and~\ref{lem:mod2-backbone}.

Thus the map
\[
(\varepsilon,r)\longmapsto\text{``implant a second $(\colC,\varepsilon,\varepsilon,\colC)$
gap at $r$''}
\]
defines a bijection from $\{\colA,\colB\}\times\widetilde{\ZZ}_n$ onto $\mathcal{A}_n$.
Hence
\[
|\mathcal{A}_n|=2(n+2).
\]
Combining this with \eqref{eq:orbit-formula} yields
\[
\zeta(G_n)=|\mathcal{L}_n|
=\frac{n}{2}\,|\mathcal{A}_n|
=\frac{n}{2}\cdot 2(n+2)
=n(n+2),
\]
as claimed.
\end{proof}

\end{document}
